\documentclass[12pt,a4paper]{article}
\usepackage[T2A]{fontenc}
\usepackage[utf8]{inputenc}
\usepackage[russian]{babel}
\usepackage{geometry}
\geometry{left=21mm,right=21mm,top=20mm,bottom=18mm}
\usepackage{amsmath,amssymb,mathtools}
\usepackage{array,booktabs,longtable}
\usepackage{xcolor}
\usepackage{tikz}
\usepackage[colorlinks=true,linkcolor=blue!55!black,urlcolor=blue!55!black]{hyperref}
\setlength{\parindent}{0pt}
\setlength{\parskip}{4pt}
\setlength{\emergencystretch}{2em}
\newcommand{\mistake}{\textcolor{red}{Ошибка:}}
\newcommand{\keyidea}{\textcolor{blue!60!black}{Ключевая идея:}}
\newcommand{\result}{\textcolor{teal!60!black}{Итог:}}
\begin{document}

\begin{titlepage}
\centering
\vspace*{0.28\textheight}
{\LARGE\bfseries Ревью домашней работы\par}
\vspace{22mm}
{\large Автор: Лёвин Артём Александрович\par}
\vspace{7mm}
{\large 20 октября 2026 года\par}
\vfill
\begin{tikzpicture}
  \draw[blue!28,line width=0.8pt] (-3.3,0) .. controls (-1.8,0.45) and (-0.7,-0.42) .. (0,0)
    .. controls (0.9,0.40) and (2,-0.34) .. (3.3,0);
\end{tikzpicture}
\vspace*{12mm}
\end{titlepage}

\newpage
\section*{Сводная таблица ошибок}
\small
\setlength{\tabcolsep}{3pt}
\renewcommand{\arraystretch}{1.35}
\begin{tabular}{@{}p{0.09\linewidth}p{0.17\linewidth}p{0.26\linewidth}p{0.20\linewidth}p{0.18\linewidth}@{}}
\toprule
\textbf{№ задания} & \textbf{Тема} & \textbf{Суть ошибки} & \textbf{Ваш ответ} & \textbf{Правильный ответ} \\
\midrule
\hyperref[task:2]{2} & Система линейных уравнений, метод Гаусса & При обратной подстановке неверно вычислена первая неизвестная. & $(-2;\,-1;\,2)$ & $(0;\,-1;\,2)$ \\
\bottomrule
\end{tabular}
\normalsize

\newpage
\thispagestyle{empty}
\vspace*{\fill}
\begin{center}
{\LARGE\bfseries Разбор ошибок}
\end{center}
\vspace*{\fill}
\newpage

\phantomsection\label{task:2}
\section*{Задание № 2}
\subsection*{Текст задания}
Решить систему линейных уравнений:
\[
\begin{cases}
 x+y+2z=3,\\
 2x-y+z=3,\\
 3x-y=1.
\end{cases}
\]

\subsection*{Ваше решение}
Вы записали расширенную матрицу системы:
\[
\left(\begin{array}{rrr|r}
1&1&2&3\\ 2&-1&1&3\\ 3&-1&0&1
\end{array}\right).
\]
Используя преобразования строк, Вы обнулили элементы под первым ведущим элементом и получили матрицу, эквивалентную следующей:
\[
\left(\begin{array}{rrr|r}
1&1&2&3\\ 0&-3&-3&-3\\ 0&-4&-6&-8
\end{array}\right).
\]
Затем из утроенной третьей строки Вы вычли учетверённую вторую строку. В результате последнее уравнение стало таким:
\[
-6z=-12.
\]
Вы верно нашли $z=2$, а из второго уравнения --- $y=-1$. Далее в работе записано:
\[
x+(-1)+4=3,\qquad \boxed{x=-2}.
\]

\subsection*{Ваш ответ}
\[
\boxed{(-2;\,-1;\,2)}.
\]

\subsection*{Ошибка}
\mistake\ арифметическая ошибка при нахождении $x$ на последнем шаге обратной подстановки. Сам метод исключения неизвестных применён правильно.

\subsection*{Где именно возникает ошибка}
Последний правильный переход в Вашем решении:
\[
x+(-1)+4=3.
\]
Первый неправильный переход:
\[
x+(-1)+4=3\quad\Longrightarrow\quad x=-2.
\]
Именно это значение $x$ не следует из записанного равенства. Значения $y=-1$ и $z=2$ найдены верно.

\subsection*{Почему этот шаг неверен}
При подстановке найденных значений двух неизвестных в первое уравнение система сводится к простому линейному уравнению с одной неизвестной:
\[
x-1+4=3.
\]
Сначала нужно объединить известные числа в левой части. Сумма минус единицы и четырёх равна трём, поэтому
\[
x+3=3.
\]
Теперь вычитаем три из обеих частей равенства:
\[
x=3-3=0.
\]
Если же взять указанное Вами значение $x=-2$, то левая часть первого исходного уравнения равна
\[
-2+(-1)+2\cdot2=-2-1+4=1,
\]
а должна равняться трём. Поэтому записанный ответ не удовлетворяет уже первому уравнению системы. Необходимо проверять не только преобразования матрицы, но и арифметику при окончательном восстановлении неизвестных.

\subsection*{Ключевая идея}
\keyidea\ после приведения расширенной матрицы к ступенчатому виду неизвестные находятся в обратном порядке. На каждом шаге следует сначала подставить уже найденные значения, аккуратно упростить числовые слагаемые и только затем выразить оставшуюся неизвестную. В конце нужно проверить всю найденную тройку в исходной системе.

\subsection*{Исправление от последнего правильного шага}
Исправлять весь метод Гаусса не требуется. Достаточно вернуться к верно записанному равенству:
\begin{align*}
x+(-1)+4&=3,\\
x+3&=3,\\
x&=3-3,\\
x&=\boxed{0}.
\end{align*}
Таким образом, при сохранении Ваших правильных значений $y=-1$ и $z=2$ получаем тройку $(0;\,-1;\,2)$.

\subsection*{Полное правильное решение}
Покажем весь ход метода Гаусса, чтобы проверить все преобразования независимо.

\textbf{Шаг 1.} Составим расширенную матрицу исходной системы:
\[
\left(\begin{array}{rrr|r}
1&1&2&3\\ 2&-1&1&3\\ 3&-1&0&1
\end{array}\right).
\]

\textbf{Шаг 2.} Вычтем из второй строки удвоенную первую строку:
\[
R_2\leftarrow R_2-2R_1.
\]
Покомпонентно получим: $2-2=0$, $-1-2=-3$, $1-4=-3$, $3-6=-3$. Поэтому
\[
\left(\begin{array}{rrr|r}
1&1&2&3\\ 0&-3&-3&-3\\ 3&-1&0&1
\end{array}\right).
\]

\textbf{Шаг 3.} Вычтем из третьей строки утроенную первую строку:
\[
R_3\leftarrow R_3-3R_1.
\]
Здесь $3-3=0$, $-1-3=-4$, $0-6=-6$, $1-9=-8$. Получаем
\[
\left(\begin{array}{rrr|r}
1&1&2&3\\ 0&-3&-3&-3\\ 0&-4&-6&-8
\end{array}\right).
\]

\textbf{Шаг 4.} Устраним неизвестную $y$ из третьей строки. Заменим третью строку разностью утроенной третьей и учетверённой второй:
\[
R_3\leftarrow 3R_3-4R_2.
\]
Действительно, коэффициент при $y$ равен $3\cdot(-4)-4\cdot(-3)=0$, при $z$ --- $3\cdot(-6)-4\cdot(-3)=-6$, а правая часть --- $3\cdot(-8)-4\cdot(-3)=-12$. Значит,
\[
\left(\begin{array}{rrr|r}
1&1&2&3\\ 0&-3&-3&-3\\ 0&0&-6&-12
\end{array}\right).
\]
Все выполненные замены строк сохраняют множество решений системы: мы прибавляли к строкам кратные других строк, а в последней операции коэффициент при изменяемой строке был ненулевым.

\textbf{Шаг 5.} Выполним обратную подстановку. Из третьей строки следует:
\[
-6z=-12\quad\Longrightarrow\quad z=\frac{-12}{-6}=2.
\]
Из второй строки:
\begin{align*}
-3y-3z&=-3,\\
-3y-3\cdot2&=-3,\\
-3y-6&=-3,\\
-3y&=3,\\
y&=-1.
\end{align*}
Из первой строки:
\begin{align*}
x+y+2z&=3,\\
x-1+2\cdot2&=3,\\
x-1+4&=3,\\
x+3&=3,\\
x&=0.
\end{align*}
Следовательно,
\[
\boxed{(x;y;z)=(0;\,-1;\,2)}.
\]

\subsection*{Проверка}
Подставим полученные числа во \emph{все три исходных уравнения}, а не только в то, из которого находили $x$.

Первое уравнение:
\[
0+(-1)+2\cdot2=-1+4=3.
\]
Второе уравнение:
\[
2\cdot0-(-1)+2=0+1+2=3.
\]
Третье уравнение:
\[
3\cdot0-(-1)=0+1=1.
\]
Во всех трёх случаях получены нужные правые части. Следовательно, найденная тройка действительно является решением исходной системы.

\subsection*{Итог}
\result\ преобразования методом Гаусса и значения $y=-1$, $z=2$ у Вас правильные. Ошибка допущена только в последнем арифметическом действии при нахождении $x$. Верный ответ:
\[
\boxed{(0;\,-1;\,2)}.
\]

\subsection*{Задача на закрепление}
Решите систему линейных уравнений методом Гаусса. После нахождения всех неизвестных проверьте полученную тройку подстановкой в каждое исходное уравнение:
\[
\begin{cases}
x+y+z=3,\\
2x-y+z=8,\\
3x-y=5.
\end{cases}
\]

\vspace{8mm}
\begin{center}
\small Лёвин Артём Александрович эксклюзивно для Анны Трапезниковой
\end{center}
\end{document}
